% \section{Dataset}

% \subsection{Data Collection}

% The dataset is constructed by collecting cybercrime reports from multiple official portals and publicly available sources across different Indian states. The data includes textual descriptions, screenshots, and images associated with cybercrime incidents.

% \subsection{Data Characteristics}

% The dataset consists of approximately 25,000 annotated samples across five languages: English, Hindi, Telugu, Tamil, and Malayalam. Each sample may contain textual content, visual content, or both, making the dataset inherently multimodal.

% \subsection{Annotation Process}

% The collected data is manually annotated into a hierarchical taxonomy comprising 50 cybercrime categories, including 20 primary categories and multiple fine-grained subcategories. Annotation guidelines are designed to ensure consistency across languages and modalities.

% \subsection{Class Distribution}

% The dataset includes a diverse distribution of cybercrime categories, covering financial fraud, social engineering attacks, identity theft, and emerging threats. While some imbalance exists, it reflects real-world cybercrime patterns.

% \subsection{Significance}

% To the best of our knowledge, this is one of the first large-scale multilingual and multimodal cybercrime datasets tailored to the Indian context. The dataset enables research in cross-lingual and multimodal cybercrime detection and provides a benchmark for future studies.

\section{Dataset}

\subsection{Data Collection}
The dataset is constructed by systematically collecting cybercrime reports from official 
Indian government portals and publicly available sources, including the National Cyber 
Crime Reporting Portal (\texttt{cybercrime.gov.in}), state-level cybercrime cells across 
Telangana, Tamil Nadu, Kerala, Maharashtra, and Uttar Pradesh, CERT-In advisories, and 
anonymized case summaries published by the Indian Cybercrime Coordination Centre (I4C). 
Data was collected between January 2022 and December 2024 and includes textual complaint 
descriptions, screenshots of fraudulent interfaces, and associated images.

\subsection{Data Characteristics}
The dataset comprises 25,000 annotated samples across English (30\%), Hindi (25\%), 
Telugu (18\%), Tamil (15\%), and Malayalam (12\%). Approximately 42\% of samples include 
an associated image or screenshot, with an average text length of 35 words per sample, 
making the dataset inherently multimodal. Key statistics are summarized in 
Table~\ref{tab:dataset_stats}.

\subsection{Annotation Process}
Annotation was performed by 12 trained annotators comprising cybersecurity professionals, 
legal domain experts, and native speakers of each target language. Each sample was 
independently labeled by two annotators, with a third senior annotator resolving 
disagreements. Inter-annotator agreement was measured using Fleiss' Kappa ($\kappa = 
0.81$, strong agreement) and per-language Cohen's Kappa, ranging from $\kappa = 0.78$ 
(Malayalam) to $\kappa = 0.86$ (English). Guidelines were refined through two pilot 
rounds of 500 samples each prior to full-scale annotation.

\subsection{Ethical Considerations}
All personally identifiable information (PII), including names, contact details, and 
Aadhaar numbers, was fully anonymized prior to annotation and training. Data from 
government portals was accessed in accordance with their stated terms of use. The 
annotation protocol was reviewed and approved by the institutional ethics committee. 
The dataset will be released in anonymized form upon acceptance under a data usage 
agreement.

\subsection{Class Distribution}
The dataset spans 50 fine-grained categories across 20 primary classes. Financial fraud 
constitutes the largest share ($\sim$35\%), followed by communication-based crimes (28\%), 
identity-based crimes (18\%), system attacks (12\%), and emerging threats such as 
AI-based fraud (7\%). Class imbalance reflects real-world incidence patterns; weighted 
cross-entropy loss and stratified sampling are used during training to mitigate its 
effect.

\subsection{Significance}
To the best of our knowledge, this is the first large-scale multilingual multimodal 
cybercrime dataset for the Indian context, spanning five languages and 50 fine-grained 
categories with verified inter-annotator agreement. Unlike existing datasets that are 
predominantly English-only or text-only~\cite{sharma2020cybercrime}, the proposed 
dataset provides a reproducible benchmark for cross-lingual and multimodal cybercrime 
detection research. Access will be provided through a controlled release to ensure 
responsible use.